% !PRINTED 390
\weberpage{390}
Betrachten wir nun irgend eine primitive ganzzahlige quadratische
 Gleichung
\weq{16}{A + B\Omega + C\Omega^2 = 0,}
in der $A$, $B$, $C$ ganze Zahlen ohne gemeinsamen Theiler sind,
die der Bedingung
\[
B^2 - 4AC = D
\]
genügen, so sind die beiden Wurzeln $\Omega$, $\Omega'$ dieser Gleichung
quadratische, zur Discriminante $D$ gehörige Irrationalzahlen, und
wenn wir sie in Kettenbrüche entwickeln, so werden wir nach
\S.\,124 endlich auf Schlusszahlen $\omega_1$, $\omega_2$ kommen, die zu den reducirten
 gehören, und die also in den oben besprochenen Perioden
enthalten sein müssen.

Es ist noch festzustellen, ob diese Schlusszahlen $\omega_1$, $\omega_2$ in
derselben oder in verschiedenen Perioden enthalten sind.

Nach 6. sind sie dann und nur dann in derselben Periode
enthalten, wenn $\Omega$ und $\Omega'$ mit einander äquivalent sind. Nennen
wir solche quadratische Irrationalzahlen, die mit ihren conjugirten
 äquivalent sind, \gesperrt{zweiseitige Zahlen}\footnote{\glqq\gesperrt{Zweiseitig}\grqq\ brauchen wir nach \gesperrt{Dedekind}'s Vorschlag statt des
\gesperrt{Gauss}'schen \gesperrt{anceps} oder des sonst üblichen \gesperrt{ambig}: \gesperrt{Dirichlet-Dedekind}, Vorlesungen über Zahlentheorie. 4. Auflage, 1894, S. 139.}, so können wir
also sagen, dass die Kettenbruchentwickelung der Wurzeln der
Gleichung (16) zu einer oder zu zwei verschiedenen Perioden
führt, je nachdem die Wurzeln zweiseitig sind oder nicht.
Welcher von beiden Fällen eintritt, kann man an der Periode
selbst erkennen.

Wenn nämlich $\omega$ eine reducirte Zahl ist, und $\omega'$ conjugirt
zu $\omega$, so ist auch $-1 : \omega'$ reducirt; und wenn $\Omega$ äquivalent ist
mit $\omega$, so ist $\Omega'$ äquivalent mit $\omega'$, also auch mit $-1 : \omega'$ (\S.\,120).

Ist
\[
\omega = \{a, b, c\},
\]
so ist nach (12):
\[
-\frac{1}{\omega'} = \{c, b, a\}.
\]

Nach \S.\,124, (9) schliesst sich also $-1 : \omega'$ ebenso an $-1 : \omega_1'$
an, wie $\omega_1$ an $\omega$. Ist daher die Periode von $\omega$
\[
\omega, \omega_1, \omega_2 \ldots \omega_{\nu-1},
\]
so ist die Periode von $-1 : \omega'$:
\[
\frac{-1}{\omega'_{\nu-1}}, \frac{-1}{\omega'_{\nu-2}} \cdots \frac{-1}{\omega_1'}, \frac{-1}{\omega'}.
\]
